% ============================================================================
% Chapter III -- Architecture Design and Security
% ============================================================================
\chapter{Architecture Design and Security}
\label{chap:architecture}

\section*{Introduction}

This chapter presents the system architecture and security design in detail,
from global topology down to individual controls.

% ----------------------------------------------------------------------------
\section{Global System Architecture}
\label{sec:global-arch}
% ----------------------------------------------------------------------------

The system follows a split-deployment model with clear separation between
frontend delivery and backend processing. \figref{fig:global-arch} summarizes
the production topology end to end: the \textbf{Browser} connects to the
\textbf{Frontend Layer (Vercel)} over HTTPS, where \textbf{Next.js~15} serves
the UI and proxies API traffic server-side; requests reach the \textbf{VPS},
where \textbf{Nginx Proxy Manager} terminates inbound connections and forwards
them into a \textbf{Docker Compose} stack running \textbf{FastAPI} and
\textbf{PostgreSQL~16}; the backend then queries \textbf{Snusbase} and free
\textbf{Enrichment APIs} (LeakCheck, EmailRep.io, Cert Spotter) over HTTPS.
No breach credentials are persisted locally---only user accounts and search
metadata are stored in PostgreSQL.

\begin{figure}[H]
\centering
\includegraphics[
  width=\linewidth,
  keepaspectratio
]{architecture-diag.png}
\caption{Global system architecture.}
\label{fig:global-arch}
\end{figure}

\textbf{Key architectural decisions:}

\begin{enumerate}[nosep]
  \item \textbf{Split deployment:} The frontend is hosted on Vercel (automatic
        HTTPS, global CDN), while the backend runs on a dedicated VPS. This
        separates static asset delivery from API processing and database I/O.
  \item \textbf{Rewrite proxy pattern:} In production, Next.js rewrites all
        \texttt{/api/v1/*} requests to the backend server-side. The browser
        never makes a direct cross-origin HTTP call, eliminating mixed-content
        issues when the API is exposed on plain HTTP behind the VPS.
  \item \textbf{Docker containerization:} The VPS backend is orchestrated with
        \texttt{docker-compose.yml}: a \textbf{PostgreSQL} container on an
        internal bridge network and a \textbf{FastAPI} container built from a
        multi-stage Alpine Dockerfile (non-root user, health checks, CPU/memory
        limits, \texttt{no-new-privileges}). The database is not reachable from
        the public internet.
  \item \textbf{Nginx Proxy Manager:} On the VPS, Nginx Proxy Manager sits in
        front of the Docker stack as the reverse proxy entry point. It routes
        incoming API traffic to the backend container, supports TLS termination
        when a custom subdomain is configured, and coexists with other projects
        on the same host through hostname-based routing.
  \item \textbf{Stateless backend:} FastAPI stores no server-side sessions.
        Authentication state is carried entirely in JWT access and refresh
        tokens, which simplifies scaling and container restarts.
  \item \textbf{External data sources:} Breach records are fetched in real
        time from Snusbase and enrichment APIs; only search audit metadata
        (query, type, result count, timestamp) is written to PostgreSQL.
\end{enumerate}

% ----------------------------------------------------------------------------
\section{Backend API Design}
\label{sec:api-design}
% ----------------------------------------------------------------------------

The backend exposes a RESTful API under the \texttt{/api/v1} prefix. All
endpoints return JSON and follow consistent error formats.

\begin{table}[ht]
\centering
\caption{REST API endpoints.}
\label{tab:endpoints}
\small
\begin{tabularx}{\textwidth}{@{} l l l X @{}}
\toprule
\textbf{Method} & \textbf{Endpoint} & \textbf{Auth} & \textbf{Description} \\
\midrule
POST & \texttt{/auth/register} & None &
  Create account (email, password, invite code) \\
POST & \texttt{/auth/login} & None &
  Authenticate; returns access + refresh tokens \\
POST & \texttt{/auth/refresh} & None &
  Exchange refresh token for new access token \\
GET  & \texttt{/auth/me} & Bearer &
  Return current user profile \\
POST & \texttt{/search} & Analyst+ &
  Search leaked databases (7 types + wildcard) \\
GET  & \texttt{/history} & Bearer &
  Return user's last 200 searches \\
GET  & \texttt{/health} & None &
  Health check (returns \texttt{\{status: ok\}}) \\
\bottomrule
\end{tabularx}
\end{table}

Each protected endpoint uses FastAPI's dependency injection system. The
\texttt{CurrentUser} dependency extracts and validates the JWT from the
\texttt{Authorization} header. The \texttt{SearchUser} dependency
additionally enforces role requirements via \texttt{require\_role(ANALYST,
ADMIN)}.

% ----------------------------------------------------------------------------
\section{Search Engine Pipeline}
\label{sec:search-pipeline}
% ----------------------------------------------------------------------------

The search pipeline orchestrates validation, API calls, enrichment, and
result aggregation concurrently. \figref{fig:search-pipeline} shows how a
query flows from the API layer through the Search Engine to Snusbase and
supplementary API Services before results are scored, deduplicated, and
returned.

\begin{figure}[H]
\centering
\includegraphics[
  width=\linewidth,
  keepaspectratio
]{search-diag.png}
\caption{Search engine pipeline.}
\label{fig:search-pipeline}
\end{figure}

\subsection{Input Validation}

The \texttt{validators} module enforces type-specific rules:

\begin{itemize}[nosep]
  \item \textbf{Email:} Must match a standard email regex.
  \item \textbf{Domain:} Must be a valid domain extracted via
        \texttt{tldextract}.
  \item \textbf{IP:} Validated against IPv4/IPv6 format.
  \item \textbf{Hash:} Must match MD5 (32~hex), SHA-1 (40~hex), or
        SHA-256 (64~hex) patterns.
  \item \textbf{Username/Name/Password:} Minimum 2 characters, maximum 320.
  \item \textbf{Wildcard mode:} Relaxes all format checks but enforces a
        3-character minimum to prevent overly broad queries.
\end{itemize}

\subsection{Concurrent Execution}

The aggregator uses a \texttt{ThreadPoolExecutor} with 4 workers to run the
Snusbase API call and enrichment calls in parallel. Since the Snusbase client
(\texttt{httpx}) uses synchronous HTTP, \texttt{asyncio.to\_thread} is not
needed---the executor offloads blocking I/O from the event loop naturally.

\subsection{Severity Scoring}

Each leak record is assigned a severity level based on the fields present:

\begin{itemize}[nosep]
  \item \textbf{Critical:} Password in plaintext or cracked hash present.
  \item \textbf{High:} Hash present (but no plaintext password).
  \item \textbf{Medium:} Email and username exposed, but no credentials.
  \item \textbf{Low:} Minimal exposure (e.g., only an email address).
\end{itemize}

% ----------------------------------------------------------------------------
\section{Security Architecture}
\label{sec:security}
% ----------------------------------------------------------------------------

This section details nine security layers forming a defense-in-depth strategy
aligned with OWASP best practices~\cite{owasp-top10}.

\subsection{Authentication: JWT Access and Refresh Tokens}

Authentication uses a dual-token scheme:

\begin{itemize}[nosep]
  \item \textbf{Access token:} Short-lived (15~minutes), signed with
        HS256 using a 256-bit secret. Contains the user ID and role in its
        claims. Sent in the \texttt{Authorization: Bearer} header.
  \item \textbf{Refresh token:} Longer-lived (7~days), used exclusively to
        obtain a new access token. Transmitted in the request body, never
        stored in cookies.
\end{itemize}

Password hashing uses \texttt{bcrypt} with 12 rounds (configurable via
\texttt{BCRYPT\_ROUNDS}), producing a 60-character hash that includes the
salt.

\begin{lstlisting}[language=Python,caption={JWT token creation (simplified).}]
def create_access_token(sub: str, extra_claims: dict) -> str:
    payload = {
        "sub": sub,
        "type": "access",
        "exp": datetime.utcnow()
              + timedelta(minutes=settings.ACCESS_TOKEN_EXPIRE_MINUTES),
        **extra_claims,  # includes {"role": "analyst"}
    }
    return jwt.encode(payload, settings.SECRET_KEY,
                      algorithm=settings.ALGORITHM)
\end{lstlisting}

\subsection{Authorization: Role-Based Access Control (RBAC)}

Access control is enforced at the API layer through FastAPI's dependency
injection:

\begin{lstlisting}[language=Python,caption={RBAC dependency injection.}]
def require_role(*roles: UserRole):
    def checker(user: User = Depends(get_current_active_user)):
        if user.role not in roles:
            raise HTTPException(403, "Insufficient permissions")
        return user
    return checker

# Applied to the search endpoint:
SearchUser = Annotated[User,
    Depends(require_role(UserRole.ANALYST, UserRole.ADMIN))]
\end{lstlisting}

Three roles are defined: \texttt{admin} (full access), \texttt{analyst}
(search + history), and \texttt{viewer} (history only, no search).

\subsection{Anti-Enumeration: Constant-Time Login}

A common attack vector is \emph{timing-based user enumeration}: an attacker
measures login response times to determine whether an email exists in the
system (bcrypt verification takes measurable time; skipping it is fast).

The system mitigates timing-based user enumeration by running a dummy bcrypt
verification when the account is not found:

\begin{lstlisting}[language=Python,caption={Constant-time login defense.}]
_DUMMY_HASH = hash_password("dummy-password-for-constant-time")

if user is None:
    verify_password(form_data.password, _DUMMY_HASH)  # same timing
    raise auth_error
\end{lstlisting}

Combined with the generic error message ``Incorrect email or password,'' this
makes it infeasible to distinguish between a non-existent account and a wrong
password.

\subsection{Registration Gate: Invite Code}

To prevent unauthorized sign-ups, registration requires a matching invite
code when \texttt{INVITE\_CODE} is configured:

\begin{lstlisting}[language=Python,caption={Invite-code verification with constant-time comparison.}]
if settings.INVITE_CODE:
    if not hmac.compare_digest(provided, settings.INVITE_CODE):
        raise HTTPException(403,
            "A valid invite code is required to register.")
\end{lstlisting}

The use of \texttt{hmac.compare\_digest} prevents timing attacks on the
invite code itself. The registration error for duplicate emails is also
generic (``Unable to register with the provided details'') to avoid
confirming which emails are already registered.

\subsection{Rate Limiting}

Per-endpoint rate limiting is implemented using \texttt{slowapi} (a
Starlette/FastAPI wrapper around \texttt{limits}):

\begin{table}[ht]
\centering
\caption{Rate limits per endpoint category.}
\label{tab:rate-limits}
\begin{tabular}{@{} l c @{}}
\toprule
\textbf{Endpoint category} & \textbf{Limit (per minute)} \\
\midrule
Authentication (\texttt{/auth/*}) & 5 \\
Search (\texttt{/search})         & 20 \\
General API                       & 60 \\
\bottomrule
\end{tabular}
\end{table}

Rate limits are keyed by client IP. When exceeded, the API returns
\texttt{429 Too Many Requests} with a \texttt{Retry-After} header.

\subsection{HTTP Security Headers}

The backend injects strict security headers on every response via a custom
middleware (\texttt{SecurityHeadersMiddleware}):

\begin{table}[ht]
\centering
\caption{HTTP security headers.}
\label{tab:headers}
\small
\begin{tabularx}{\textwidth}{@{} l X @{}}
\toprule
\textbf{Header} & \textbf{Value and Purpose} \\
\midrule
\texttt{X-Content-Type-Options} &
  \texttt{nosniff} --- prevents MIME-type sniffing attacks \\
\texttt{X-Frame-Options} &
  \texttt{DENY} --- blocks clickjacking via iframes \\
\texttt{Referrer-Policy} &
  \texttt{strict-origin-when-cross-origin} --- limits referrer leakage \\
\texttt{Permissions-Policy} &
  Disables geolocation, microphone, camera, payment, USB \\
\texttt{Cross-Origin-Opener-Policy} &
  \texttt{same-origin} --- isolates browsing context \\
\texttt{Strict-Transport-Security} &
  \texttt{max-age=63072000; includeSubDomains; preload} (production only) \\
\bottomrule
\end{tabularx}
\end{table}

The frontend additionally sets a \textbf{Content Security Policy (CSP)} via
\texttt{next.config.js}, restricting script sources, connection targets, and
frame ancestors.

\subsection{CORS: Strict Origin Allowlist}

Cross-Origin Resource Sharing is configured with an explicit allowlist of
origins (set via the \texttt{CORS\_ORIGINS} environment variable). In
production, only the Vercel frontend domain is permitted. Credentials
(\texttt{allow\_credentials=True}) and specific HTTP methods are explicitly
listed. Wildcard origins (\texttt{*}) are never used.

\subsection{Input Validation}

All user input passes through Pydantic models (FastAPI's built-in validation
layer) before reaching business logic:

\begin{itemize}[nosep]
  \item Email fields validated by \texttt{pydantic[email-validator]}.
  \item Password: minimum 12 characters, strength-checked (uppercase,
        lowercase, digit, special character).
  \item Search terms: type-specific regex, maximum 25 terms per query,
        320-character limit per term.
  \item Invite code: maximum 128 characters.
\end{itemize}

\subsection{Docker Security}

The containerization strategy follows Docker security best practices:

\begin{enumerate}[nosep]
  \item \textbf{Alpine base image:} Minimal attack surface (5~MB base vs.
        120~MB for Debian).
  \item \textbf{Multi-stage build:} Build dependencies (compilers, headers)
        are excluded from the production image.
  \item \textbf{Non-root execution:} The application runs as user
        \texttt{leaklens} (UID~1000), never as root.
  \item \textbf{\texttt{no-new-privileges}:} The Docker security option
        prevents privilege escalation inside the container.
  \item \textbf{\texttt{forwarded-allow-ips}:} Uvicorn trusts
        \texttt{X-Forwarded-*} headers only from the Docker bridge network
        (\texttt{127.0.0.1,172.16.0.0/12}), preventing IP spoofing from
        external clients.
  \item \textbf{\texttt{.dockerignore}:} Excludes \texttt{.env}, tests, and
        Python cache from the build context.
  \item \textbf{Resource limits:} Docker Compose enforces CPU and memory
        limits (1.5~CPU / 512~MB for backend, 1.0~CPU / 512~MB for
        PostgreSQL).
  \item \textbf{Trivy scanning:} A \texttt{trivy-scan.sh} script builds the
        image and scans it for known vulnerabilities before deployment.
\end{enumerate}

\begin{lstlisting}[language={},caption={Dockerfile excerpt --- non-root user and restricted forwarding.},label={lst:dockerfile}]
# Runtime stage
FROM python:3.12-alpine AS runtime
RUN adduser -D -u 1000 leaklens
COPY --from=builder /install /usr/local
COPY ./app /home/leaklens/app
USER leaklens
CMD ["uvicorn", "app.main:app", "--host", "0.0.0.0", "--port", "8000", \
     "--proxy-headers", "--forwarded-allow-ips=127.0.0.1,172.16.0.0/12"]
\end{lstlisting}

% ----------------------------------------------------------------------------
\section{Containerization with Docker}
\label{sec:docker}
% ----------------------------------------------------------------------------

The backend Dockerfile uses a two-stage build strategy:

\begin{figure}[ht]
\centering
\tikzfig{%
\begin{tikzpicture}[
  stage/.style={draw,rounded corners,minimum width=3.4cm,minimum height=1.1cm,
    align=center,font=\footnotesize},
  arrow/.style={-{Stealth[length=2.5mm]},thick},
]
  \node[stage,fill=blue!10] (builder) {\textbf{Builder}\\python:3.12-alpine\\pip install};
  \node[stage,fill=green!10,right=2cm of builder] (runtime)
    {\textbf{Runtime}\\python:3.12-alpine\\non-root user};
  \draw[arrow] (builder) -- node[above,font=\tiny]{COPY} (runtime);
\end{tikzpicture}%
}
\caption{Docker multi-stage build.}
\label{fig:docker-stages}
\end{figure}

The builder stage installs all Python dependencies (including C extensions
like \texttt{bcrypt} and \texttt{asyncpg}) with their build-time headers. The
runtime stage copies only the compiled packages, resulting in a lean
production image.

% ----------------------------------------------------------------------------
\section{Deployment Architecture}
\label{sec:deployment}
% ----------------------------------------------------------------------------

\begin{figure}[ht]
\centering
\tikzfig{%
\begin{tikzpicture}[
  node distance=0.9cm and 1.5cm,
  box/.style={draw,rounded corners,minimum width=2.6cm,minimum height=0.85cm,
    align=center,font=\footnotesize},
  arrow/.style={-{Stealth[length=2.5mm]},thick},
]
  \node[box,fill=gray!10] (github) {GitHub};
  \node[box,fill=green!10,below left=1.2cm and 0.2cm of github] (vercel)
    {Vercel};
  \node[box,fill=orange!10,below right=1.2cm and 0.2cm of github] (vps)
    {VPS\\Docker};
  \node[box,fill=purple!10,below=0.8cm of vps] (npm) {Nginx PM};
  \draw[arrow] (github) -- (vercel);
  \draw[arrow] (github) -- (vps);
  \draw[arrow] (npm) -- (vps);
\end{tikzpicture}%
}
\caption{Deployment workflow.}
\label{fig:deployment}
\end{figure}

\begin{itemize}[nosep]
  \item \textbf{Frontend (Vercel):} Every push to \texttt{main} triggers an
        automatic build and deployment. Vercel provides HTTPS, a global CDN,
        and preview deployments for pull requests.
  \item \textbf{Backend (VPS):} The operator SSHs into the VPS, pulls the
        latest code, and runs \texttt{docker compose up -d -{}-build}. The
        PostgreSQL data is persisted on a named Docker volume.
  \item \textbf{Nginx Proxy Manager:} Already present on the VPS for other
        projects, it routes traffic by subdomain with automatic Let's Encrypt
        certificates. The VPS can host multiple projects simultaneously, each
        with its own subdomain.
  \item \textbf{Environment variables:} Production secrets (\texttt{SECRET\_KEY},
        \texttt{POSTGRES\_PASSWORD}, \texttt{SNUSBASE\_API\_KEY},
        \texttt{INVITE\_CODE}) are stored in a \texttt{.env} file on the VPS,
        never committed to Git. A \texttt{.env.example} template is provided
        in the repository.
\end{itemize}

% ----------------------------------------------------------------------------
\section*{Conclusion}

This chapter presented the system architecture and its nine security layers.
The search pipeline handles concurrent API orchestration with input validation
and severity scoring. The next chapter shows the concrete realization through
screenshots and production metrics.
